\section{辽西、海上与朝廷：1621--1627}

\eventanchor{ML-1621-901}{天启元年后金攻取辽阳，明军退守辽西}；\eventanchor{ML-1622-901}{次年广宁又失}。辽东并非地图上一条整齐后移的线，而是城池、屯田、商路、军户与撤退人群共同构成的地带。北京收到的是败报、请饷与问责；一座城中的人则要决定逃离、留守、服役或重接交易关系。\eventanchor{ML-1622-902}{同年荷兰舰队进攻澳门而被守军击退}，又使珠江口的贸易、防务、传教士与地方官处于同一片海面。澳门议事会、耶稣会文字、明方奏疏与VOC材料能确认攻守发生，不能由港口个案推出整个东南海贸的规模。

\eventanchor{ML-1623-901}{天启三年孙承宗在宁远等辽西城堡布置防御与屯田}，试图将守备安放在可补给、可驻守的地点。明实录和奏疏可以确证朝廷决策，满文与朝鲜材料能校对对方行动；它们不授予任一方战报以最后的兵数和控制范围。\eventanchor{ML-1624-901}{天启四年明荷在澎湖交涉后，荷兰人撤往台湾南部设据点}，显示海防结果也可能是移动而非消失。福建官员面对海岛、季风和补给，VOC在寻找可长期停泊并接入贸易的地点；撤离澎湖改变水道条件，不能讲成明廷对整片海域的彻底恢复。

北京的宫门亦在重排。\eventanchor{ML-1624-902}{杨涟上疏弹劾魏忠贤等内廷势力}；\eventanchor{ML-1625-901}{次年魏忠贤集团推动逐出、处分东林相关官员}。奏疏保存作者提出的指控和程序位置，实录与后出的列传记录朝廷选择；它们不足以把“东林”或“阉党”化为内部一致、动机透明的两群人。\eventref{ML-1625-902}{同年大秦景教流行中国碑在西安被发现}，拓本、题跋和书信说明碑石被怎样发现、抄录和解释，而不直接复原唐代信仰。

\eventanchor{ML-1626-901}{天启六年袁崇焕守宁远，后金军撤退}。宁远的守与退是辽西防务重新取得的节点，却不能以一场战事抹去辽阳、广宁失守后的给养困境。\eventanchor{ML-1627-901}{翌年熹宗死，信王朱由检即位为崇祯帝}；\eventanchor{ML-1627-902}{冬季魏忠贤失势、自尽，党羽陆续被处置}。继承和处分改变北京的官僚关系，并未使辽东、沿海和各地征解暂停。

\section{没有实录的朝廷，分散的压力（1628--1633）}

\eventanchor{ML-1628-901}{崇祯元年王二等在陕北起事}，\eventref{ML-1628-902}{陕西、山西旱灾、蝗灾与赈济记录也在累积}。不能将自然灾害、催征、武装流动与每一位参与者的动机归为一个原因；驿夫、马匹、铺舍和沿线市场所承受的节省，会改变离乡者、地方官和武装集团可作的选择。崇祯朝没有存世《明实录》，此后《崇祯长编》、奏疏、地方志、朝鲜与清方材料只能拼出谨慎的先后线，不能填成逐日透明的朝廷日记。

\eventanchor{ML-1629-901}{崇祯二年后金军经长城北段入关、进逼北京}。长城不能自行阻断军事移动：关口、道路、守兵、粮秣和情报的缺口须分别追问。\eventref{ML-1629-902}{同年徐光启等组织历法修订}，表明朝廷仍能集结技术与文书资源；项目设立不能证明知识在各地同样使用。\eventanchor{ML-1630-901}{崇祯三年袁崇焕被处决}，后出的《明史》与清初叙事常将其编入成败判断；《崇祯长编》、相关奏疏和清方材料可以比对程序与说法，仍不能凭一条叙事裁定密谋与责任。\eventref{ML-1630-902}{同年在辽饷之外议定剿饷}，财政压力开始以新的名称进入征解链。

\eventanchor{ML-1631-901}{崇祯四年孔有德部在吴桥起兵}，欠饷与地方冲突相互交织；\eventanchor{ML-1631-902}{大凌河城被围后降}，辽东局势再坏。\eventanchor{ML-1632-901}{次年孔有德、耿仲明部攻陷登州}，登莱海防和军政秩序受冲击。\eventref{ML-1632-902}{剿饷征解形成催征与缓征的行政链}，定额、欠额、实解和家庭负担不能互换。\eventanchor{ML-1633-901}{崇祯六年郑芝龙率明军在料罗湾击败荷兰舰队}；福建奏疏与VOC档案相互校验时，首先显示不同政权和商业组织怎样描述船只、港口与胜负，而不能令“海盗”“商人”“明军”成为不变身份。\eventref{ML-1633-902}{徐光启同年去世}，其农政、历法项目由门人续加整理，文本延续不等于制度效果自动延续。

\section{多省边界、朝鲜与三饷（1634--1640）}

\eventanchor{ML-1634-901}{崇祯七年张献忠一度受抚}，部众去向和约束随即成为湖广地方军政问题。\eventanchor{ML-1635-901}{次年流动作战集团在多省边界重组，明军分区追剿}；受抚、转徙与追剿不是相同的行政结果。\eventanchor{ML-1635-902}{郑芝龙同年受命清剿刘香等海上武装}，也表明沿海贸易和海防秩序须在港口、人手与地方协作中重新安排。

\eventanchor{ML-1636-901}{崇祯九年皇太极改国号为清}；\eventanchor{ML-1636-902}{同年清军再次入侵朝鲜，朝鲜被迫改变对清关系}。北京、盛京与汉城的材料有不同历法、外交语言和战争立场，辽东战场并不止于明清边线。\eventref{ML-1637-901}{翌年朝廷增设练饷}，三饷并行的征解压力加重；\eventanchor{ML-1638-901}{崇祯十一年清军再度入关，京畿及北方多地受劫掠}，防务和征解的压力继续向腹地传递。\eventref{ML-1640-901}{崇祯十三年西北、中原多地旱蝗和粮食危机持续}，奏疏、地方志与灾异记录能定位何地进入朝廷视野，不能相加为无误的全国死亡或饥民数字。

\section{城门、河道与关中基地（1641--1644）}

\eventanchor{ML-1641-901}{崇祯十四年李自成攻克洛阳，福王朱常洵被杀}；\eventanchor{ML-1641-902}{张献忠同年攻取襄阳，襄王朱翊铭被杀}。王府所在城市成为争夺粮食、军政据点和政治名义的场所，每个城内的暴力和居民遭遇仍须由地方材料逐项核验。

\eventanchor{ML-1642-901}{崇祯十五年松山、锦州战局失败，洪承畴降清}，辽西防御再受重创；\eventanchor{ML-1642-902}{九月黄河决口后开封受灾}。前者来自明、清双方军政记录，后者来自河渠叙述、奏疏与开封地方材料。将二者同称“崩溃”会遮蔽不同空间条件：一个是边防与军队困局，一个是河道、城池与居民面对水患。

\eventanchor{ML-1643-901}{崇祯十六年李自成取襄阳，以新顺王名义整编势力}；\eventanchor{ML-1643-902}{十一月攻取西安，将关中作为大顺政权基地}。\eventanchor{ML-1644-901}{翌年正月在西安建国号大顺并称帝}。这些步骤显示政治组织、城市据点和军队调动的先后，不能以后来进入北京的结果抹平其中仍存在的谈判、供给和区域差异。

\eventref{ML-1644-902}{崇祯十七年三月十九日，李自成军进入北京，崇祯帝死于煤山。}本卷的时间线在此止步。北京地方材料、《崇祯长编》末卷、《明史·思宗本纪》与《甲申传信录》能够交叉定位这一结局，却带有后出编纂、幸存者记忆或政治解释。南明与清军入关后的战争和政权重组，属于下一卷，不能反向补写进这一日的叙事。

\begin{debate}[没有崇祯实录，年序怎样仍可成立]
崇祯朝无存世《实录》。《崇祯长编》与奏疏保存行政时间，陕西、山西、河南、湖广、福建和开封地方材料限定区域经验；满文、《清太宗实录》《仁祖实录》、VOC档案及近时见闻提供不同观察位置。它们足以支持事件顺序和行动者，不能消除军数、伤亡、动机、命令传达和社会损失的争议。
\end{debate}